\clearpage

\lehead[]{\normalfont\sffamily\hspace*{-2.00cm}\textcolor{white}{\colorbox{lightblue}{\makebox[1.60cm][r]{\thechapter}}}\hspace{0.17cm}\textcolor{lightblue}{\chaptertitle}}
\rohead[]{\textcolor{lightblue}{\chaptertitle}\normalfont\sffamily\hspace*{0.17cm}\textcolor{white}{\colorbox{lightblue}{\makebox[1.60cm][l]{\thechapter}}}\hspace{-2.00cm}}
%\chead[]{}
\rehead[]{\textcolor{lightblue}{AvHG, Inf, My}}
\lohead[]{\textcolor{lightblue}{AvHG, Inf, My}}

\lstset{style=myJava}

\section{Arbeiten mit Zeichenketten (Strings)}

%Im Zusammenhang mit Benutzeroberflächen muss regelmäßig mit Zeichenketten
%gearbeitet werden.

Im Package \myPackage{java.lang} wird die Klasse \myClass{String} zur Arbeit mit
Zeichenketten definiert. Das Package \myPackage{java.lang} wird immer
automatisch eingebunden und braucht deshalb nicht explizit importiert zu
werden.

Mit Objekten der Klasse \myClass{String} haben wir schon oft gearbeitet. Jede in
doppelte Anführungszeichen gesetzte Zeichenkette (zum Beispiel \texttt{"Hallo"}) ist ein
Objekt der Klasse \myClass{String}. Im Gegensatz zu anderen Klassen brauchen
Objekte der Klasse \myClass{String} nicht mit \lstinline|new| erzeugt zu
werden. Wie bei \lstinline|int|- und \lstinline|boolean|-Variablen kann man
einer \myClass{String}-Variable einfach einen Wert durch ein
Gleichheitszeichen zuweisen:

\begin{lstlisting}
String s1 = "Hallo";
String s2 = "Guten Tag";
\end{lstlisting}

Mit dem Plus-Zeichen können zwei Strings aneinandergehängt werden:

\begin{lstlisting} 
s1 = s1 + s2;           æ// in s1 steht jetzt: "HalloGuten Tag"; s2 bleibt unverändert 
æs1 = s2 + " Otto";      æ// in s1 steht jetzt: "Guten Tag Otto"
æs1 += "!";              æ// in s1 steht jetzt: "Guten Tag Otto!"
\end{lstlisting}

\begin{framed}
Der Nachteil bei der Arbeit mit Strings ist, dass Strings in Java unveränderlich
(\emph{immutable}) sind: Bei jeder Veränderung einer Zeichenkette wird ein neues
Objekt im Speicher erzeugt. Wenn beispielsweise \lstinline|s1| einen neuen
Inhalt bekommt, bleibt der alte Inhalt im Speicher \glqq als Leiche\grqq\
liegen, bis irgendwann der sogenannte \emph{Garbage Collector}
(\glqq Müll-Aufsammler\grqq ) der Java Virtual Machine Zeit hat, diesen
Speicher wieder freizugeben. Wenn man eine effizientere Speichernutzung haben
will, muss man mit der Klasse \myClass{StringBuilder} arbeiten. Das lohnt sich
aber nur, wenn man ein Programm schreibt, in dem große Mengen von
Zeichenketten verwaltet bzw.\ verändert werden müssen.
\end{framed}

\vfill

\subsection{Methoden der Klasse \myClass{String} (Auswahl)}

\noindent \emph{Hinweis:} Genau wie bei Arrays beginnt die Zählung der Positionen
der einzelnen Zeichen immer bei \texttt{0}.

\medskip
\begin{minipage}{1.0\textwidth} % Seitenumbruch verhindern

\bgroup
\def\arraystretch{1.2}
\begin{tabularx}{\textwidth}{|p{85mm}|X|}
\hline
\texttt{int length()} & 
Länge des Strings (Anzahl der Zeichen)
\\ \hline
\texttt{String toLowerCase()}

\texttt{String toUpperCase()} & 
Umwandlung in Klein- bzw.\ Großschreibung
\\ \hline
\texttt{char charAt(int index)} & 
liefert das Zeichen an Position \texttt{index}
\\ \hline
\texttt{String replace(char oldChar, char newChar)} & 
ersetzt überall im String das Zeichen \texttt{oldChar} durch \texttt{newChar}
(funktioniert auch mit ganzen Teilstrings)
\\ \hline
\texttt{String trim()} & 
entfernt alle Leerzeichen am Anfang und Ende
\\ \hline
\texttt{int indexOf(int ch)} 

\texttt{int indexOf(String str)} & 
Anfangsposition des Zeichens oder Teilstrings. Falls das Zeichen nicht im
String vorkommt, wird \texttt{-1} zurückgegeben.
\\ \hline
\texttt{String substring(int beginIndex)} 

\texttt{String substring(int beginIndex, int endIndex)} & 
liefert den Teilstring von Position \texttt{beginIndex} bis eine Position
vor \texttt{endIndex} (oder bis zum Ende)
\\ \hline
\texttt{boolean contains(String str)} & 
\texttt{true}, falls \texttt{str} im gegebenen String-Objekt enthalten ist;
sonst \texttt{false}
\\ \hline
\end{tabularx}
\egroup

\subsubsection{Vergleich von Strings}

\bgroup
\def\arraystretch{1.2}
\begin{tabularx}{\textwidth}{|p{85mm}|X|}
\hline
\texttt{==} & 
vergleicht die Speicheradressen (Referenzen) der Objekte
\\ \hline
\texttt{boolean equals(Object anObject)} & 
vergleicht den Inhalt der Zeichenketten 
(überschriebene Methode der Superklasse \myClass{Object})
\\ \hline
\texttt{boolean equalsIgnoreCase(String anotherString)} & 
vergleicht Zeichenketten ohne Beachtung von Groß- und Kleinschreibung
\\ \hline
\texttt{int compareTo(String anotherString)} & 
Lexikografischer Vergleich:

\texttt{this < anotherString}: negativer Wert

\texttt{this == anotherString}: $0$

\texttt{this > anotherString}: positiver Wert
\\ \hline
\end{tabularx}
\egroup

\end{minipage} % Seitenumbruch verhindern


\section{Methoden der Klasse \myClass{Character}}

Die Klasse \myClass{Character} besitzt folgende statische Methoden, die für
dich nützlich sind:

\medskip
\bgroup
\def\arraystretch{1.2}
\begin{tabularx}{\textwidth}{|p{85mm}|X|}
\hline
\texttt{public static boolean isDigit(char ch)} & 
Ziffer von '0' bis '9'?
\\ \hline
\texttt{public static boolean isLetter(char ch)} & 
Buchstabe?
\\ \hline
\texttt{public static boolean isUpperCase(char ch)} & 
Großbuchstabe?
\\ \hline
\texttt{public static boolean isLowerCase(char ch)} & 
Kleinbuchstabe?
\\ \hline
\end{tabularx}
\egroup
